\section{The Italian Legal System}
\label{sec:domains}
In this section, we describe the Italian legal system, the different types of trials with the documents involved, and the relationship between Italian law and \acl{ai}.

The Italian Republic uses the \emph{civil law} system, in contrast with countries such as the United Kingdom of Great Britain and Northern Ireland and the United State of America, which use the \emph{common law} system. Civil law is based on codes, while common law is based on case-law~\cite{Carlson2009}. 
Typically, civil law systems rely on four foundational codes~\cite{Carlson2009,apple1995primer}:
\begin{itemize}
    \item \emph{Civil code}: governs private relations among individuals and organizations;
    \item \emph{Civil procedure code}: sets the rules for how civil disputes are adjudicated;
    \item \emph{Criminal code}: defines crimes and the corresponding penalties;
    \item \emph{Criminal procedure code}: regulates the investigation, prosecution, and adjudication of criminal offenses.
\end{itemize}

These codes are intended as coherent and consistent principles for supporting the broader legal system. In contrast, common law systems typically develop codes in response to specific issues, often drafting them only after principles have already emerged through case law~\cite{Carlson2009}.
Furthermore, in common law, precedent cases are of central importance, as judges are required to follow earlier decisions, and a single ruling from a higher court may determine the outcome of a case. This is different in civil law, where a single precedent case is generally not treated as binding, with the possible exception of decisions in the constitutional sphere~\cite{Carlson2009}. 

The organizational structure of civil law court systems is broadly similar to that of common law systems, featuring a court of first instance, a court of appeals, and a supreme court. Additionally, specialized courts exist with jurisdiction limited to particular areas of law. Unlike in common law jurisdictions, juries are generally not used in civil law systems. The cases are decided by judges, eventually including lay judges~\cite{Carlson2009}, such as in the Italian court of assize~\cite[Art.~3]{legge10apr1951}---which is competent to judge serious criminal offenses, such as the ones carrying maximum penalties of at least 24 years of imprisonment~\cite[Art.~5]{CPP}.  

In the Italian system, both in civil and in criminal procedures, the three grades of jurisdiction operate as follows. The courts of first instance decide the facts and merits of a case in their entirety~\cite[Arts. 99--310]{CPC}; \cite[Arts. 429--544]{CPP}. The courts of appeal then provide a second, comprehensive review to ensure both factual and legal correctness~\cite[Arts. 339--359]{CPC}; \cite[Arts. 593--605]{CPP}, thereby giving effect to the constitutional guarantee of challenging judicial decisions~\cite[Art.~111]{Costituzione}. Finally, the Court of Cassation represents the third and highest grade, serving as the supreme authority tasked not with re-examining facts but with ensuring the uniform interpretation and application of the law, as regulated in the civil~\cite[Art.~360--382]{CPC} and criminal~\cite[Art.~606--613]{CPP} codes of procedure.
Decisions of the Court of Cassation are not laws and are not formally binding, yet they strongly influence lower courts and promote a consistent interpretation of the law, reflecting the court's so-called \emph{nomophilactic function}~\cite{lazovic2022}. In Italy, the supreme court has a dual role: to render a final judgment on individual cases, thereby ensuring justice, and to establish interpretative guidelines that shape national jurisprudence~\cite[Art.~65]{RD1941_12_art65}.


\subsection{Legal Proceedings}

In the Italian legal system, a legal proceeding may arise in either the civil or the criminal sphere, with distinct initiation mechanisms reflecting the nature of the dispute.

% \subsubsection{Civil Proceedings}

Civil proceedings are initiated by a private party. The \emph{plaintiff} begins the process by filing a writ of summons with the competent court. This document must indicate: the identities of the parties, the tribunal being addressed, the factual and legal basis of the claim, the evidence intended to be produced, and the date of the hearing \cite[Art.~163]{CPC}. 

The \emph{defendant} responds with a statement of defense, outlining their defenses and any counterclaims \cite[Art.~167]{CPC}. From this moment, the judicial authority schedules hearings, manages the collection of evidence, and ultimately renders a decision---a \emph{judgment}---which must include motivations as required by the Constitution~\cite[Art.~111]{Costituzione}.

% \subsubsection{Criminal Proceedings}

\medskip
\noindent
Criminal proceedings begin with the occurrence of a potential criminal offense. When the report of a crime is received, the public prosecutor opens a preliminary investigation~\cite[Art.~326]{CPP}.

During this phase, investigative activities are performed by the public prosecutor and the judicial police, and documented in official records~\cite[Art.~373]{CPP}. The suspect is formally notified through a guarantee notice, a document informing them of the charges and their rights~\cite[Art.~369]{CPP}. If sufficient evidence exists, the prosecutor requests an indictment and the judge for the preliminary hearing decides whether the case proceeds to trial \cite[Art.~416]{CPP}.



% \subsubsection{Documents}
\medskip
\noindent
The documentation of civil and criminal proceedings is governed by the provisions of the respective codes of procedure~\cite[Art.~168--169]{CPC}; \cite[Art.~431, 114]{CPP}, which govern the formation, content, and accessibility of case files.
Civil files typically include the summons, payment receipts, pleadings, briefs, hearing minutes, judicial orders, procedural documents related to the taking of evidence, and judgments~\cite[Art.~168]{CPC}, while criminal trial files contain procedural records related to prosecutability and civil claims, non-repeatable investigative activities documented by the police or the parties, records of evidentiary proceedings, official documents, and, when applicable, seized items~\cite[Art.~431]{CPP}.
Furthermore, the content of the judgment are regulated, requiring identification of the parties, a statement of facts, the legal reasoning, and the operative part of the decision~\cite[Art.~546]{CPP}; \cite[Art.~132]{CPC}. For this reason, judgments can be treated as semi-structured documents, in the sense that they usually contain identifiable sections, even though the expression and organization of legal language may vary---for example, depending on the style of the judge~\cite{juliahb2024,galli2025analisi,pont2023legalsummarisationllmsprodigit}.



% \subsection{Analyzing Legal Documents}

% improving access to legal information, facilitating search and retrieval of relevant case law, generating summaries and keywords, extracting arguments, and even offering statistical forecasts of outcomes. These functions promise to enhance efficiency in overloaded court systems, reduce inconsistencies in jurisprudence, and support judges and lawyers in handling the growing volume and complexity of legal data. Importantly, AI is conceived not as a substitute for judicial decision-making, but as an auxiliary instrument that can strengthen transparency, consistency, and access to justice~\cite{juliahb2024}.
% ADELE project and outcome prediction

% chat
% **Why AI can be useful in justice**
% Artificial intelligence (AI) is increasingly seen as a valuable support tool for judicial systems. According to the JuLIA Handbook, AI can assist in a variety of tasks: improving access to legal information, facilitating search and retrieval of relevant case law, generating summaries and keywords, extracting arguments, and even offering statistical forecasts of outcomes. These functions promise to enhance efficiency in overloaded court systems, reduce inconsistencies in jurisprudence, and support judges and lawyers in handling the growing volume and complexity of legal data. Importantly, AI is conceived not as a substitute for judicial decision-making, but as an auxiliary instrument that can strengthen transparency, consistency, and access to justice.

% **What has been done so far in Italy**
% Italy has begun experimenting with AI tools, especially in the context of predictive justice. The Documentation Centre (CED) of the Court of Cassation has explored algorithmic methods for citation analysis and legal research, while research projects such as ADELE have tested applications for summarization, argument extraction, and outcome prediction. The 2023 Cartabia Reform introduced a shift from “documents” to “data,” laying the groundwork for more systematic use of AI in judicial administration. Legislative attention has also increased: Italy presented a proposal (Bill no. 1146/2024) to regulate AI specifically in the justice sector, aligning with the EU Artificial Intelligence Act but adopting a cautious approach that limits AI's role to supporting research and organizational tasks.

% **Issues and regulatory safeguards**
% The JuLIA Handbook underscores that AI in justice is a high-risk domain, requiring strong safeguards. Risks include algorithmic opacity, biased datasets, and potential erosion of judicial independence if predictive tools were given too much weight. For this reason, both European and Italian frameworks emphasize principles of transparency, non-discrimination, due process, and human oversight. The Italian Council of State has further clarified that algorithmic systems must comply with constitutional requirements of impartiality and accountability. Overall, Italy's approach reflects a careful balance: acknowledging AI's potential to make justice more efficient and accessible, while embedding safeguards to protect judicial autonomy and fundamental rights.
%%%

% https://www.wired.it/article/ddl-ai-senato-approvazione-regole-controllo-autorita/
% https://www.repubblica.it/tecnologia/2025/09/18/news/legge_italiana_intelligenza_artificiale_deep_fake-424854661/
\subsection{Artificial Intelligence in the Legal Domain}
\label{sec:background:ai4legal}

\ac{ai} offers considerable opportunities for the legal sector, fundamentally reshaping traditional workflows through the automation of routine tasks and the enhancement of complex decision-making processes \cite{juliahb2024}. The application of \ac{ai} is not intended to supplant human legal expertise but rather to act as a sophisticated form of decision-support systems (DSSs), augmenting the capabilities of legal professionals. Such systems can perform an array of tasks, from conducting advanced precedent searches to predicting potential case outcomes, thereby improving efficiency while maintaining the integrity of human judicial oversight \cite{juliahb2024}.

An example of a practical \ac{ai} tool for the legal domain is the ADELE project~\cite{galli2024analytics}, which addressed different legal domains, including Value added tax (VAT) and trademarks and patents. The project aimed to support legal research and decision-making through several functions. These included the automatic extraction of citations between legal documents and the construction of citation networks to identify influential cases and clusters of similar decisions. ADELE also applied deep learning models to extract summaries and keywords from legal texts, facilitating information retrieval and organization. In addition, the project experimented with argument extraction and with predicting case outcomes by examining possible correlations between the arguments advanced by the parties and the decisions of the courts~\cite{juliahb2024}.

For this thesis, the most important use case for \ac{ai} in the legal domain is automatically analyzing and organizing---\eg by criteria such as relevance to a specific case---vast quantities of documents, avoiding time-consuming and expensive manual operations. Indeed, \ac{ai} systems have been applied in common law jurisdiction to identify relevant documents among millions~\cite{juliahb2024}. However, also in civil law jurisdiction, some cases may require the analysis of a high number of related decisions to ensure the uniform interpretation of the law~\cite{egea2024does}. Furthermore, criminal investigations frequently involve the examination of vast amounts of heterogeneous documents, such as web pages, chat records, investigative reports, or expert analyses~\cite{szekely2015building,real_batini2021semantic,pozzi_wise_2025}.

The use of \ac{ai}-based systems in this context enables users to satisfy information needs through multiple modalities, such as faceted search~\cite{kejriwal2017searchengine}, which allows filtering based on entities mentioned in documents, \acl{qa}~\cite{hoppe_2021_legalqa_german,sovrano2020legal}, and \acl{sta}. The latter can support reading by highlighting salient elements in a judgment~\cite{swoopes2025interfacedesignsupportlegal,ziwei_abstractexplorer_2025,raymond_skimming_sci_2023} (\eg sections or parties). These use cases motivate part of the work conducted during my \phd, in particular the research presented in the following publications, which I co-authored: \textcite{pozzi_wise_2025}, \textcite{BELLANDI2024}, \textcite{aixia_2023}, and in \Cref{chap:extraction}.



In the context of information retrieval and organization, the importance of data annotation cannot be overstated. Projects in Italy, such as the PNRR-PON ``Giustizia Agile'' (Italian for ``agile justice''), have focused on creating labeled datasets for Italian legal documents to support further \ac{ai} research. This included the crucial task of \acf{ner} (described in \Cref{sec:information_extraction}), which is foundational for enabling \ac{ai} systems to accurately identify and classify key entities within legal texts, such as persons, organizations, or mentioned laws and regulations, which can be used to improve the retrieval mechanisms \cite{juliahb2024}, \eg with faceted search~\cite{ARMENTANO2014_nlp_faceted,hearst_2006_faceted}.

\subsection{Artificial Intelligence Regulations}

The integration of artificial intelligence (\ac{ai}) into the legal and judicial
system presents significant challenges, particularly concerning fundamental
rights, transparency, and data protection. The European Union has addressed
these concerns through the \emph{Artificial Intelligence Act (Regulation (EU)
2024/1689)}~\cite{ai_act}, often referred to as \emph{\acs{ai} Act}, which classifies \ac{ai} systems used in the
administration of justice as high-risk~\cite[Annex~III]{ai_act}. While the
\ac{ai} Act does not prohibit their use, it mandates \emph{human
oversight}---\ie that the final decision-making authority remains with a
human~\cite[Art.~14]{ai_act}---ensuring that \ac{ai} functions as a support
mechanism rather than an autonomous decision-maker~\cite{juliahb2024}. This
aligns with the European Ethical Charter on the Use of \ac{ai} in Judicial
Systems~\cite{ethicalcharter2018}, which emphasizes principles such as
transparency, non-discrimination, and due process.

These principles have been further consolidated in Italy with the approval of the Law on Artificial Intelligence (L.~1146/2024)~\cite{leggeia2025}. This national law integrates and strengthens the European framework, reaffirming the anthropocentric approach whereby \ac{ai} in the judicial sector may be used exclusively for instrumental and support activities. It also introduces criminal provisions, including penalties for the illicit dissemination of \ac{ai}-generated or manipulated content, thereby mitigating emerging risks~\cite{reuters_italy_ai_law_2025}.

These regulatory developments are already reflected in judicial decisions across both Italy and the EU. The Court of Bologna, for instance, has considered cases involving algorithmic decision-making, focusing on principles of equality and non-discrimination. Similarly, the Court of Justice of the European Union (CJEU), in the \emph{SCHUFA Holding} case~\cite{schufaholding2023}, examined whether an \ac{ai}-based credit scoring system constitutes automated decision-making under the GDPR~\cite{gdpr2016}.
Together, these examples highlight the ongoing legal and ethical discourse on the governance of \ac{ai} in justice systems.

A further crucial aspect concerns \emph{privacy} and the lawful processing of data used to train or operate \acp{llm}. Legal and judicial documents often contain personal data~\cite[Art.~4]{gdpr2016}, which must be processed lawfully, fairly, and for specified purposes in accordance with the \emph{General Data Protection Regulation (GDPR, Regulation (EU) 2016/679)}~\cite[Art.~5]{gdpr2016}. In this context, data controllers---the entities determining the purposes and means of data processing~\cite[Art.~4.7]{gdpr2016}---are encouraged to adopt \emph{anonymization} or \emph{pseudonymization} techniques~\cite[Art.~4.5]{gdpr2016} to reduce privacy risks, as these measures limit or remove the identifiability of individuals while still allowing useful data analysis.
% However, only anonymization---where re-identification is impossible---fully exempts data from the GDPR’s scope, whereas pseudonymized data remain subject to its provisions. 
In particular, the controller must ensure that data are not transferred outside the European Economic Area unless an adequate level of protection is guaranteed~\cite[Art.~44--49]{gdpr2016}. 
The use of \ac{api}-based models complicates compliance with these principles, as such systems typically involve transferring data to external servers operated by third-party providers, thereby making it more difficult to ensure adherence to the principles of data minimization and purpose limitation~\cite[Art.~5.1.c--b]{gdpr2016}.




Under the \ac{ai} Act~\cite{ai_act}, the \emph{provider} is the entity that develops or markets the \ac{ai} system---for example, a company offering a model through an online \ac{api}~\cite[Art.~3.3]{ai_act}---whereas the \emph{deployer} is the natural or legal person, public authority, or other body that uses the system under its authority~\cite[Art.~3.4]{ai_act}. In legal settings, the deployer may correspond to a court or law enforcement agency integrating an \ac{ai} model to support document review or investigative analysis. When the model is accessed through an external \ac{api}, these entities must establish a lawful processing basis (\eg consent or public interest)~\cite[Art.~6 and~9]{gdpr2016}, define roles and safeguards via a data processing agreement~\cite[Art.~28.3]{gdpr2016}, and comply with additionally transfer rules if data is processed outside the EU~\cite[Art.~44--49]{gdpr2016}. 

Moreover, the \ac{ai} Act reinforces these obligations for high-risk systems, requiring providers and deployers to implement risk management, data governance, and logging mechanisms for traceability~\cite[Art.~9--12]{ai_act}, and to ensure that the system preserves human oversight and data security~\cite[Art.~14--15]{ai_act}. When using remote \acp{api}, these safeguards depend on the provider’s infrastructure, making compliance difficult to demonstrate. Conversely, locally deployed models allow institutions to retain control over the data processing environment, enforce internal access policies, and ensure that confidential information remains within the organization's secure perimeter---a crucial requirement for judicial and investigative applications.

\subsection{Illustrative Examples of Domain Data}

To conclude the section, \Cref{fig:civil-judgment-legal-articles,fig:chat-investigation-example} present two illustrative examples: a civil judgment and an investigative chat log extracted from a seized suspect's smartphone.





\begin{figure}[htbp]
\centering
\begin{minipage}{0.9\linewidth}
\small
\textbf{ORDINARY COURT OF BOLOGNA}\\
Civil Division\\
Judgment No.~987/2025\\
Published on February~14,~2025\\[0.5em]
\textbf{ITALIAN REPUBLIC}\\
IN THE NAME OF THE ITALIAN PEOPLE\\[0.5em]
The Court, sitting as a single judge, in the person of Judge Laura Conti, has delivered the following:\\[0.5em]
\textbf{JUDGMENT}\\[0.3em]
in the civil case registered under No.~3210/2024\\
between\\[0.3em]
Giovanni Bianchi, residing in Bologna, represented and defended by Attorney Paolo De~Santis, with elected domicile at his office in Bologna --- plaintiff\\[0.3em]
and\\[0.3em]
Davide Ricci, residing in Modena, represented and defended by Attorney Silvia Ferri, with elected domicile at her office in Modena --- defendant\\[0.5em]

\textbf{FACTS}\\
The plaintiff summoned the defendant to court, alleging that on April~3,~2022, he lent him the sum of EUR~10,000, as evidenced by a private written agreement signed by both parties, with repayment due by December~31,~2022, pursuant to Art.~1813~c.c. Despite repeated requests, the defendant failed to repay the sum. The defendant claimed that he had already reimbursed part of the debt in cash and that the remaining amount had been compensated by the plaintiff's use of his vehicle, pursuant to Art.~1241~c.c. (compensation of debts).\\[0.5em]

\textbf{REASONS FOR THE DECISION}\\
The documentary evidence and the bank records produced confirm the existence of a loan contract under Art.~1813~c.c. and disprove any effective compensation under Art.~1241~c.c. No evidence was provided of partial repayment. Pursuant to Art.~2697~c.c., the burden of proof rests on the party alleging payment. The plaintiff's version is consistent and corroborated by witness statements under Arts.~2721--2725~c.c.\\
Therefore, the defendant must return the sum of EUR~10,000, plus statutory interest under Art.~1282~c.c. from January~1,~2023 until payment.\\[0.5em]

\textbf{RULING}\\
The Court, definitively ruling, dismissing all other claims,\\
orders Mr.~Davide~Ricci to pay Mr.~Giovanni~Bianchi EUR~10,000 plus statutory interest from January~1,~2023 until payment in full (Arts.~1282 and~1284~c.c.);\\
orders the defendant to reimburse the plaintiff's legal expenses, quantified at EUR~1,800 plus VAT and surcharges, pursuant to Arts.~91--92~c.p.c.\\[0.5em]
Bologna, February~10,~2025\\
\textit{The Judge}\\
Laura~Conti
\end{minipage}
\caption[Illustrative example of an Italian civil judgment.]{Illustrative example of an Italian civil judgment. c.c. stands for Italian Civil Code~\cite{CC} and c.p.c. for Italian Code of Civil Procedure~\cite{CPC}.}
\label{fig:civil-judgment-legal-articles}
\end{figure}





\begin{figure}[htbp]
\centering
\begin{minipage}{0.92\linewidth}
\small
\textbf{Participants:} Giovanni Bianchi, Davide Ricci\\
\textbf{Date Range:} 2025--04--22 to 2025--04--25\\[0.3em]

\textbf{[2025/04/22 -- 11:32]} Giovanni Bianchi: Davide, we need to talk.\\
\textbf{[2025/04/22 -- 11:35]} Davide Ricci: About what now?\\
\textbf{[2025/04/22 -- 11:37]} Giovanni Bianchi: You know what. You ignored the court order again.\\
\textbf{[2025/04/22 -- 11:39]} Davide Ricci: I told you, I don't have that money anymore. Things went bad with the workshop.\\
\textbf{[2025/04/22 -- 11:40]} Giovanni Bianchi: Don't lie. I saw your bank statements, the transfer went straight to that trip to Milan and that ``business dinner'' at \textit{Hotel Duomo}.\\
\textbf{[2025/04/22 -- 11:42]} Davide Ricci: It was all work-related.\\
\textbf{[2025/04/22 -- 11:45]} Giovanni Bianchi: Meet me Friday, 19:00, at Bar Centrale. Bring the papers.\\
\textbf{[2025/04/22 -- 11:46]} Davide Ricci: Fine.\\[0.3em]

\textbf{[2025/04/24 -- 18:51]} Giovanni Bianchi: \textit{[Voice message --- 2~min 58~sec]}\\
\quad \textit{Transcription:}\\
\quad Listen, Davide, I'm done playing games. The Court of Bologna, Judge Conti herself, said you must pay me back --- it's all there in the papers. You keep saying you're broke, but I know about that account at Banca Emilia Romagna and your new tool deal with Officina Rinaldi in Modena. You spent my money on that weekend in Milan, dinner at Grand Hotel et de Milan and even those tickets for the race at the Monza Circuit.\\
\quad I checked --- you still have that Piaggio van parked behind your shop in Via Emilia. Don't tell me you sold it. You've been moving money through your cousin's account at Credito Padano.\\
\quad If you don't bring something --- even half --- on Friday, I swear, Davide, I'll find you myself. I'll be waiting near the pharmacy next to Piazza Maggiore, around six-thirty. Don't make me do something stupid.\\[0.3em]

\textbf{[2025/04/24 -- 19:02]} Davide Ricci: Calm down. I'll come, okay?\\
\textbf{[2025/04/24 -- 19:03]} Giovanni Bianchi: We'll see.\\
\end{minipage}
\caption{Illustrative example of chat exchange.}
\label{fig:chat-investigation-example}
\end{figure}



%%%%%%%%%%%%%%%%%%%%%%%%%%%%%%%%%%%%%%%%%%%%%%
% documents for the picture
% \section*{Example Case: Cross-Document Evidence Integration}

% \subsection*{Civil Judgment}

% \noindent

% \bigskip
% \subsection*{WhatsApp Conversation}

% \noindent


% \bigskip
% \subsection*{Bank Transfer Summary}

% \begin{center}
% \begin{tabulary}{\linewidth}{LLLLLL}
% \toprule
% \textbf{Date} & \textbf{Sender Account} & \textbf{Receiver Account} & \textbf{Amount (EUR)} & \textbf{Description} & \textbf{Notes} \\
% \midrule
% 2022/04/03 & IT12X0306909601 & IT34R0306909602 & 10,000 & Private loan agreement & From Bianchi to Ricci \\
% 2022/05/12 & IT34R0306909602 & IT55A0306909603 & 2,300 & ``Hotel Duomo Milano S.p.A.'' & Trip to Milan \\
% 2022/06/01 & IT34R0306909602 & IT78L0306909604 & 1,800 & ``Officina Rinaldi SRL'' & Workshop expenses \\
% 2023/09/15 & IT34R0306909602 & --- & --- & --- & Account closed (insolvent) \\
% \bottomrule
% \end{tabulary}
% \end{center}

% \bigskip
% \subsection*{Police Report (Extract)}

% \noindent
% \textbf{Police Headquarters of Bologna}\\
% Report No.~452/B/2025\\
% Filed on April~26,~2025\\[0.3em]
% On April~25,~2025, at approximately 19:15, local police units intervened near \textbf{Via Indipendenza~102, Bologna}, following reports of an altercation. Upon arrival, officers found \textbf{Davide Ricci}, 41~years old, resident of Modena, lying on the pavement with severe head trauma. Emergency services declared him deceased at the scene.\\[0.3em]
% Eyewitnesses reported seeing \textbf{Giovanni Bianchi}, 45~years old, leaving the nearby alley shortly after a heated argument. CCTV footage from a pharmacy camera shows both individuals meeting at approximately 18:37 and moving toward the adjacent parking area.\\[0.3em]
% The investigation established that the meeting coincided with the conversation exchanged through WhatsApp on April~24,~2025.\\
% An arrest warrant was issued for \textbf{Giovanni Bianchi} on suspicion of homicide (Art.~575~c.p.).\\[0.5em]
% Bologna, April~26,~2025\\
% \textit{Inspector}\\
% Elena Rossi
